% Conversão de chapter-06-text-h-0122; origem imutável RELATORIO_ATUAL.odt.
\hypertarget{__RefHeading___Toc4855_3118176633}{}
\refstepcounter{section}
\section*{5.4 Conclusão}
\label{ages-iv:text-h-0122}
\addcontentsline{toc}{section}{5.4 Conclusão}

% editorial-recorda-conclusion-p01
\indent Durante muito tempo, quis atuar como AGES IV e me perguntei se estaria preparado para a responsabilidade que via nessa função. No Recorda, entendi o que essa responsabilidade significava no dia a dia. Eu precisava decidir o que priorizar, acompanhar o trabalho da equipe e lidar com as consequências das escolhas que fazíamos. Foi nesse processo, ao tomar decisões, reconhecer erros e rever minha forma de trabalhar, que fui desenvolvendo o preparo que eu esperava ter antes de assumir o papel.\par

% editorial-recorda-conclusion-p02
\indent Aprendi a dar à equipe condições para trabalhar com mais autonomia. Isso envolveu detalhar as tarefas, identificar dependências e delegar atividades de acordo com as prioridades de cada Sprint. Também precisei reconhecer quando um bloqueio exigia meu apoio e quando os integrantes podiam seguir por conta própria. As retrospectivas nos ajudaram a fazer esse ajuste ao longo do semestre. O que discutíamos depois de uma entrega se refletia na maneira de organizar as squads e preparar as atividades da Sprint seguinte.\par

% editorial-recorda-conclusion-p03
\indent Essa experiência também mudou minha forma de pensar sobre estratégia. Alterar a prioridade de uma tarefa ou redistribuir atividades pode influenciar o ritmo de desenvolvimento, a qualidade da entrega e a maneira como as pessoas trabalham juntas. Discutir essas escolhas com os demais AGES IV me ajudou a considerar os efeitos de cada decisão sobre o projeto inteiro. Assumir a responsabilidade pelas decisões também significou responder pelo que entregamos e pelo que ficou pendente, reconhecer meus próprios erros e atuar junto à equipe para corrigir os problemas que surgiam.\par

% editorial-recorda-conclusion-p04
\indent Considero que tive uma atuação muito satisfatória no projeto. Ver o time resolver com mais independência dificuldades que antes exigiam nossa intervenção foi uma das partes mais gratificantes da minha experiência no Recorda. Essa evolução me ajudou a compreender o valor do trabalho de gestão em uma entrega construída por toda a equipe.\par

